\subsection{线性映射的有理性}

定义 3. 设 \(\mathbf{V}_1, \mathbf{V}_2\) 是 \(\mathbf{K}\) 上的两个右向量空间，分别带有 \(\mathbf{K}'\) -结构 \(\mathbf{V}_1'\) , \(\mathbf{V}_2'\) 。一个 K-线性映射 \(f: \mathbf{V}_1 \to \mathbf{V}_2\) 称为在 \(\mathbf{K}'\) 上有理的，如果 \(f(\mathbf{V}_1') \subset \mathbf{V}_2'\) 。

若 \(V_{3}\) 是 K 上的第三个右向量空间，带有一个 \(K^{\prime}\) -结构 \(V_{3}^{\prime}\) ，且 K-线性映射 \(g:V_{2}\rightarrow V_{3}\) 在 \(K^{\prime}\) 上有理，则显然 \(g\circ f:V_{1}\rightarrow V_{3}\) 在 \(K^{\prime}\) 上有理。

命题 3. 设 \(\mathbf{V}_1, \mathbf{V}_2\) 是 \(\mathbf{K}\) 上的两个右向量空间，且 \(\mathbf{V}_1', \mathbf{V}_2' \mathbf{K}'\) -结构分别给定在 \(\mathbf{V}_1, \mathbf{V}_2\) 上。 \(\mathbf{V}_1\) （相应地 \(\mathbf{V}_2\) ）典范地等同于 \(\mathbf{V}_1' \otimes_{\mathbb{K}'} \mathbf{K}\) （相应地 \(\mathbf{V}_2' \otimes_{\mathbb{K}'} \mathbf{K}\) ）（第 1 号，命题 1）。

\begin{enumerate}
\def\labelenumi{(\roman{enumi})}
\item
  映射 \(f' \mapsto f' \otimes 1_{\mathbf{K}} = f_{(\mathbf{K})}'\) 是从 \(\operatorname{Hom}_{\mathbf{K}'}(\mathbf{V}_1', \mathbf{V}_2')\) 到 \(\mathbf{V}_1\) 到 \(\mathbf{V}_2\) 中的在 \(\mathbf{K}'\) 上有理的 K-线性映射集合上的一个双射；逆双射把每个 K-线性映射 \(f: \mathbf{V}_1 \to \mathbf{V}_2\) （在 \(\mathbf{K}'\) 上有理的）对应到 K'-线性映射 \(f': \mathbf{V}_1' \to \mathbf{V}_2'\) ，它与 f 到 \(\mathbf{V}_1'\) 上的限制具有相同的图像。
\item
  对每个 K-线性映射 \(f: \mathbf{V}_1 \to \mathbf{V}_2\) （在 \(\mathbf{K}'\) 上有理的），
\end{enumerate}

\[
f (\mathrm{V} _ {1} ^ {\prime}) = f (\mathrm{V} _ {1}) \cap \mathrm{V} _ {2} ^ {\prime} \quad and \quad \bar {f} ^ {- 1} (\mathrm{V} _ {2} ^ {\prime}) = \mathrm{V} _ {1} ^ {\prime} + \operatorname{Ker} (f).
\]

\begin{enumerate}
\def\labelenumi{(\roman{enumi})}
\item
  显然，在上述等同下，若 \(f': \mathbf{V}_1' \to \mathbf{V}_2'\) 是一个 K'-线性映射，则 \(f_{(\mathbf{K})}' = f' \otimes 1_{\mathbf{K}}\) 在 \(\mathbf{K}'\) 上有理，且 \(f'\) 是与 \(f_{(\mathbf{K})}'\) 到 \(\mathbf{V}_1'\) 上的限制具有相同图像的映射。反之，若 \(f: \mathbf{V}_1 \to \mathbf{V}_2\) 是一个在 \(\mathbf{K}'\) 上有理的 K-线性映射，且 \(f': \mathbf{V}_1' \to \mathbf{V}_2'\) 与 \(f\) 到 \(\mathbf{V}_1'\) 上的限制具有相同的图像，则 \(f\) 和 \(f_{(\mathbf{K})}'\) 在 \(\mathbf{V}_1'\) 上一致，后者是 \(\mathbf{V}_1\) 在 \(\mathbf{K}\) 上的一个生成系，故 \(f = f_{(\mathbf{K})}'\) 。
\item
  若 \(f = f' \otimes 1_{\mathbb{K}}\) ，则 \(f(\mathbf{V}_1) = f(\mathbf{V}_1' \otimes_{\mathbb{K}'} \mathbf{K}) = f'(\mathbf{V}_1') \otimes_{\mathbb{K}'} \mathbf{K}\) ，且由于 \(f'(\mathbf{V}_1') \subset \mathbf{V}_2'\) ，有 \(f(\mathbf{V}_1') = f'(\mathbf{V}_1') = f(\mathbf{V}_1) \cap \mathbf{V}_2'\) （§ 7，第 9 号，命题 19）；公式 \(f^{-1}(\mathbf{V}_2') = \mathbf{V}_1' + \operatorname{Ker}(f)\) 立即得出。
\end{enumerate}

推论 1. 在命题 3 的记号下， \(\operatorname{Im}(f) = (\operatorname{Im}(f'))_{(\mathbb{K})}\) ，

\[
\operatorname{Ker} (f) = (\operatorname{Ker} (f ^ {\prime})) _ {(\mathbb {K})}, \quad \operatorname{Coker} (f) = (\operatorname{Coker} (f ^ {\prime})) _ {(\mathbb {K})}.
\]

特别地，要使 \(f\) 是单射（相应地，满射、零映射），必要且充分的条件是 \(f'\) 如此。若 \(f\) 是双射的，则其逆映射在 \(\mathbf{K}'\) 上有理。

这是 § 7 第 9 号命题 18 的推论的一个特例。

推论 2. 设 \(f: \mathbf{V}_1 \to \mathbf{V}_2\) 是一个在 \(\mathbf{K}'\) 上有理的 K-线性映射。对每个向量子 K-空间 \(\mathbf{W}_1\) （ \(\mathbf{V}_1\) 的；相应地 \(\mathbf{W}_2\) ， \(\mathbf{V}_2\) 的），若它在 \(\mathbf{K}'\) 上有理，则 \(f(\mathbf{W}_1)\) （相应地 \(f(\mathbf{W}_2)\) ）是 \(\mathbf{V}_2\) （相应地 \(\mathbf{V}_1\) ）的在 \(\mathbf{K}'\) 上有理的向量子 K-空间。

在命题 3 的记号下，对每个向量子 K'-空间 W \(_{1}\) （属于 V \(_{1}\) ），有 f \(_{(K)}\) (W \(_{1}\) ' ⊗ \(_{K'}\) K) = f'(W \(_{1}\) ) ⊗ \(_{K'}\) K；由此得出关于 W \(_{1}\) 的断言（第 2 号，命题 2）。另一方面，设 W \(_{2}\) 是一个向量子 K'-空间（ V \(_{2}\) 的），并令 g' 为典范 K'-线性映射 V \(_{2}\) → V \(_{2}\) /W \(_{2}\) ；于是 \(f'^{-1}(W_{2}') = \text{Ker}(g' \circ f')\) ；因此，由推论 1，

\[
\bar {f} _ {(\mathbf {K})} ^ {\prime - 1} (\mathbf {W} _ {2} ^ {\prime} \otimes_ {\mathbf {K} ^ {\prime}} \mathbf {K}) = \bar {f} ^ {\prime - 1} (\mathbf {W} _ {2} ^ {\prime}) \otimes_ {\mathbf {K} ^ {\prime}} \mathbf {K},
\]

由此得出关于 \(W_{2}\) 的断言。

设 \(V_{1}, V_{2}\) 是两个右向量 K-空间，分别带有 \(K'\) -结构 \(V_{1}', V_{2}'\) 。显然， \(V_{1}' \times V_{2}'\) 是一个 \(K'\) -结构（在 \(V_{1} \times V_{2}\) 上的），称为 \(K'\) -结构 \(V_{1}'\) 与 \(V_{2}'\) 的积。

命题 4. 要使 K-线性映射 \(f: \mathbf{V}_1 \to \mathbf{V}_2\) 在 \(\mathbf{K}'\) 上有理，必要且充分的条件是它的图像 \(\Gamma\) 在 \(\mathbf{K}'\) 上有理（对于 \(\mathbf{K}'\) -积结构，即 \(\mathbf{V}_1 \times \mathbf{V}_2\) 上的那个而言）。

设 \(g\) 为映射 \(x_1 \mapsto (x_1, f(x_1))\) （从 \(V_1\) 到 \(V_1 \times V_2\) 中的）；这是一个使得 \(\Gamma = g(V_1)\) 的 K-线性映射；若 \(f\) 在 \(K'\) 上有理，则根据命题 3 的推论 2， \(\Gamma\) 亦然。反之，设 \(\Gamma\) 在 \(K'\) 上有理，并令它具有一个 \(K'\) -结构（由 \(V_1 \times V_2\) 上的结构诱导的）；由定义立即得出，限制 \(p_1, p_2\) （到 \(\Gamma\) 上的、投影 \(pr_1, pr_2\) 的）是 K-线性映射，在 \(K'\) 上有理，分别从 \(\Gamma\) 到 \(V_1\) 和 \(V_2\) 中。由于 \(p_1\) 是双射的，其逆映射 \(q_1\) 在 \(K'\) 上有理（命题 3 的推论 1），因而 \(f = p_2 \circ q_1\) 亦然。

